% VoxShield — MICCAI 2025 Anonymous Submission
% Based on Springer LNCS template (llncs class)
\documentclass[runningheads]{llncs}
%
\usepackage[T1]{fontenc}
\usepackage{graphicx}
\usepackage{amsmath}
\usepackage{amssymb}
\usepackage{amsfonts}
\usepackage{bm}
\usepackage{booktabs}
\usepackage{multirow}
\usepackage{color}
%
% ── Uncomment for blue hyperlinks in eBook style ──
%\usepackage{hyperref}
%\renewcommand\UrlFont{\color{blue}\rmfamily}
%
\begin{document}
%
\title{VoxShield: Protecting 3D Medical Datasets from Unauthorized Training via Frequency-Aware Inter-Slice Disruption}
\titlerunning{VoxShield}
%
\author{Paper ID: 3104}

\institute{No Institute Given}
%
\maketitle
%
% ═══════════════════════════════════════════════════════════════════
%                          ABSTRACT
% ═══════════════════════════════════════════════════════════════════
\begin{abstract}
The widespread dissemination of public 3D medical image segmentation (MIS) datasets accelerates clinical research but simultaneously heightens risks of unauthorized AI model training. While \textbf{Unlearnable Examples (UE)} offer protection by injecting imperceptible perturbations to prevent effective model learning, existing methods primarily target 2D image classification. They neglect the \textit{anisotropy} and \textit{inter-slice anatomical consistency} inherent in 3D medical volumes, which are critical learning priors exploited by 3D segmentation networks. To bridge this gap, we propose \textbf{VoxShield}, a UE generation framework specifically designed for 3D MIS. VoxShield employs a UNet-based generator to produce ROI-constrained adaptive perturbations, guided by a frequency-domain adversarial loss scheme. We introduce an \textit{Inter-Slice Frequency Consistency Disruption} module that disrupts inter-slice consistency by maximizing spectral discrepancies between adjacent slices, and a \textit{Semantic Prediction Disruption} module that maximizes the $\ell_1$ divergence between clean and perturbed logits to corrupt semantic representations. Experiments on BraTS19 and FLARE21 demonstrate that VoxShield degrades 3D segmentation performance from 88\% to 7\% DSC with minimal perturbation ($\epsilon{=}4/255$), while preserving high visual fidelity.

\keywords{Unlearnable Examples \and 3D Medical Image Segmentation \and Frequency-domain Adversarial Loss \and Data Privacy Protection}
\end{abstract}
%
% ═══════════════════════════════════════════════════════════════════
%                       1  INTRODUCTION
% ═══════════════════════════════════════════════════════════════════
\section{Introduction}
\begin{figure}[t]
\centering
\includegraphics[width=\linewidth]{fig/motivation.pdf}
\caption{Motivation of VoxShield. \textbf{Left:} A model trained on clean data learns coherent volumetric representations and segments accurately on clean test volumes. \textbf{Right:} VoxShield disrupts inter-slice consistency by injecting spectrally diverse perturbations across adjacent slices; a model trained on such protected data fails to aggregate cross-slice features, producing severely degraded segmentation on clean test inputs.}
\label{fig:motivation}
\end{figure}

3D medical image segmentation (MIS) is fundamental to disease diagnosis, treatment planning, and surgical navigation. The success of modern 3D segmentation models~\cite{cciccek20163d,milletari2016v,isensee2021nnu,hatamizadeh2022unetr,tang2022swinunetr} critically depends on large-scale, voxel-annotated datasets whose creation requires substantial expert effort. Institutions increasingly share datasets through public repositories such as TCIA~\cite{tcia}, with accompanying publications indexed on PubMed~\cite{pubmed}, to facilitate collaborative research~\cite{menze2014multimodal,ma2021abdomenct}, typically under data-use agreements that restrict usage to non-commercial research purposes so as to protect patient privacy and data ownership.
 However, once published, volumetric data can be freely downloaded and repurposed for unauthorized model training. Such misuse violates both the original data-use agreements and privacy regulations such as the GDPR~\cite{gdpr2016}, which mandates that data collected for specific purposes shall not be repurposed without explicit consent~\cite{kaissis2021end}.

To counter such risks, Unlearnable Examples (UEs)~\cite{huang2021unlearnable} offer an appealing proactive defense by injecting imperceptible perturbations into training data \textit{before} release. Models trained on the perturbed data learn ``shortcut'' patterns instead of genuine semantic features, and consequently fail to generalize to clean test inputs. Representative methods include EM~\cite{huang2021unlearnable}, which crafts error-minimizing noise via bi-level optimization; TAP~\cite{fowl2021adversarial}, which leverages targeted adversarial perturbations; SEP~\cite{liu2023stable}, which stabilizes noise through a stable error-minimizing objective; LSP~\cite{yu2022availability}, which aligns perturbations with linearly separable features; REM~\cite{fu2022robust}, which enhances robustness against adversarial training; and CUDA~\cite{sadasivan2023cuda}, which generates class-wise perturbations via convolution. While UEs have been explored for 2D image classification and more recently extended to 2D segmentation~\cite{sun2024unseg} and 3D point clouds~\cite{wang2024unlearnable3d}, existing methods remain ineffective for 3D medical volumetric segmentation due to two fundamental gaps:

Existing 2D UE methods neglect the volumetric priors unique to 3D data. Modern 3D segmentation architectures rely on a key inductive bias: \textit{anatomical inter-slice consistency}~\cite{isensee2021nnu,dong2022mnet}. Organs and lesions vary smoothly along the $z$-axis, and 3D convolution kernels explicitly aggregate cross-slice context to exploit this continuity. The importance of this prior is evidenced by a well-established empirical finding: under severe anisotropy, where large inter-slice spacing naturally disrupts volumetric continuity, 2D models consistently outperform their 3D counterparts~\cite{isensee2021nnu,dong2022mnet}. Although voxel spacing is an intrinsic acquisition parameter that cannot be altered post hoc, this observation reveals that inter-slice consistency is the decisive factor governing 3D segmentation performance. This further suggests that deliberately disrupting the cross-slice aggregation process of 3D convolutions---by injecting inter-slice incoherent perturbations along the $z$-axis---can directly undermine the very inductive bias these models depend on, thereby providing a principled attack vector against 3D architectures.


As illustrated in Fig.~\ref{fig:motivation}, our key insight is that inter-slice consistency is simultaneously the source of 3D segmentation power and its exploitable vulnerability. Adjacent slices of a volume naturally share similar in-plane spatial patterns---a property that 3D convolutions heavily rely on to aggregate contextual cues along~$z$. We exploit this by injecting perturbations whose \emph{in-plane spectral characteristics} vary maximally between neighboring slices, so that 3D kernels spanning the $z$-axis receive conflicting frequency patterns at every spatial position, fundamentally undermining their ability to learn coherent volumetric representations.

Building on this insight, we propose \textbf{VoxShield}, a UE framework
tailored for 3D medical image segmentation. VoxShield comprises two core
modules: (i)~an \textit{Inter-Slice Frequency Consistency Disruption}
module that maximizes spectral discrepancies between adjacent slices in
the frequency domain, injecting ``inter-slice flickering'' that directly
undermines the $z$-axis continuity assumption of 3D models; and (ii)~a
\textit{Semantic Prediction Disruption} module that maximizes the
$\ell_1$ distance between clean and perturbed logits, forcing
perturbations to propagate through the entire network and corrupt its
semantic understanding.

% ═══════════════════════════════════════════════════════════════════
%                       2  METHODOLOGY
% ═══════════════════════════════════════════════════════════════════
\section{Methodology}

%-------------------------------------------------------------
\subsection{Preliminary}
\label{sec:preliminary}

\textbf{3D Medical Image Segmentation.}\quad
Given a clean training set
$\mathcal{D}_{\mathrm{clean}}\!=\!\{(\bm{x}_i,\bm{y}_i)\}_{i=1}^{N}$,
where $\bm{x}_i\!\in\!\mathbb{R}^{D\times H\times W}$ is a 3D medical
volume with $D$ slices and
$\bm{y}_i\!\in\!\{0,1\}^{D\times H\times W\times K}$ is the
corresponding voxel-wise annotation over $K$ classes,
3D medical image segmentation aims to optimize a network
$\mathcal{F}(\cdot\,;\theta)$ to establish the volumetric mapping from
$\bm{x}$ to $\bm{y}$.
This optimization can be formulated as:
\begin{equation}
  \arg\min_{\theta}\;
  \mathbb{E}_{(\bm{x},\bm{y})\sim\mathcal{D}_{\mathrm{clean}}}
  \bigl[\mathcal{L}_{\mathrm{seg}}\!\bigl(
    \mathcal{F}(\bm{x};\theta),\,\bm{y}\bigr)\bigr],
  \label{eq:seg}
\end{equation}
where $\mathcal{L}_{\mathrm{seg}}$ is the segmentation loss
(\textit{e.g.}, cross-entropy loss and Dice loss)
and $\theta$ denotes the trainable parameters of $\mathcal{F}$.

\textbf{Unlearnable Examples.}\quad
UEs are designed to prevent unauthorized models from extracting useful
knowledge by adding imperceptible perturbations to the training
images of $\mathcal{D}_{\mathrm{clean}}$.
For better understanding, we take the representative error-minimizing
(EM) noise~\cite{huang2021unlearnable} as an example.
EM jointly optimizes a surrogate model and sample-wise perturbations
to generate protective noise:
\begin{equation}
  \min_{\theta}\;
  \mathbb{E}_{(\bm{x},\bm{y})\sim\mathcal{D}_{\mathrm{clean}}}
  \Bigl[\min_{\bm{\delta}\in\mathcal{I}}
  \mathcal{L}_{\mathrm{seg}}\!\bigl(
    \mathcal{F}_s(\bm{x}+\bm{\delta};\theta),\,\bm{y}
  \bigr)\Bigr],
  \label{eq:ue}
\end{equation}
where $\mathcal{F}_s$ is a surrogate model,
$\mathcal{I}\!=\!\{\bm{\delta}:\|\bm{\delta}\|_\infty\!\le\!\epsilon\}$
is the feasible region that ensures imperceptibility, and
$\bm{x}+\bm{\delta}$ denotes the protected image (also called the
unlearnable example).
In this process, $\theta$ and $\bm{\delta}$ are alternately updated
to minimize $\mathcal{L}_{\mathrm{seg}}$, yielding the protected
dataset
$\mathcal{D}_{\mathrm{protect}}\!=\!\{(\bm{x}_i+\bm{\delta}_i,\,
\bm{y}_i)\}_{i=1}^{N}$.


\textbf{Task Objective.}\quad
The \textit{data protector} has access to
$\mathcal{D}_{\mathrm{clean}}$ and publishes
$\mathcal{D}_{\mathrm{protect}}$. The \textit{data exploiter} trains
a 3D segmentation architecture on
$\mathcal{D}_{\mathrm{protect}}$.
We evaluate the protection effectiveness by testing the model trained
by the exploiter on held-out clean data: if the test-time performance
degrades significantly, the protection is considered successful.
The protector aims to minimize such test-time performance while
keeping perturbations imperceptible.

%-------------------------------------------------------------
\subsection{VoxShield Framework}
\label{sec:VoxShield}
\begin{figure}[t]
  \centering
  \includegraphics[width=\textwidth]{fig/Method.pdf}
\caption{Overview of the proposed VoxShield framework. A learnable noise
generator $G$ produces mask-guided perturbations that are added to clean
3D medical images to create protected volumes. A frozen surrogate model
$F_s$ computes $\mathcal{L}_{\text{SPD}}$ and $\mathcal{L}_{\text{seg}}^{'}$
to drive protected images away from the original decision boundaries.
An inter-slice frequency consistency disruption module enforces
$\mathcal{L}_{\text{ISC}}$ by maximizing spectral discrepancy between
adjacent slices, destroying the volumetric continuity that 3D models
rely on.}
  \label{fig:method}
\end{figure}

\textbf{Overview.}\quad
Existing UEs fail in 3D medical segmentation because they ignore
\textit{inter-slice anatomical consistency}, the key inductive bias
exploited by 3D models. As illustrated in Fig.~\ref{fig:method},
VoxShield introduces two dedicated modules: (i)~\textit{Inter-Slice Frequency Consistency Disruption} maximizes the spectral
discrepancy of perturbations between adjacent slices in the frequency
domain, breaking the inter-slice continuity prior that 3D models rely on;
(ii)~\textit{Semantic Prediction Disruption} maximizes the $\ell_1$
distance between clean and perturbed logits to corrupt semantic
predictions.
To focus perturbations on anatomically meaningful regions, the raw
generator output is masked by a binary region-of-interest mask
$\bm{M}\!\in\!\{0,1\}^{D\times H\times W}$ derived from $\bm{y}$:
\begin{equation}
  \bm{\delta} = \operatorname{clamp}\!\bigl(
G(\bm{x})\odot\bm{M},\;-\epsilon,\;\epsilon\bigr),
  \label{eq:mask}
\end{equation}
where $\odot$ denotes element-wise multiplication.
The protected volume is then $\bm{x}_p\!=\!\bm{x}+\bm{\delta}$.

%------- L_ISC -------
\subsubsection{Inter-Slice Frequency Consistency Disruption.}
3D segmentation networks fundamentally assume \textit{volumetric
continuity}, i.e., anatomical structures change smoothly along $z$.
3D convolution kernels with receptive fields spanning $k_z$ slices
aggregate cross-slice features under this assumption.
Our strategy is to violate this assumption by forcing the
perturbation's spectral characteristics to vary \textit{maximally}
between consecutive slices.

For each slice $z$, we compute the 2D amplitude
spectrum of the perturbation:
$\bm{S}_z=|\mathcal{F}_{2\text{D}}(\bm{\delta}_z)|\in\mathbb{R}^{H\times W}$,
where $\bm{\delta}_z\!\in\!\mathbb{R}^{H\times W}$ is the
$z$-th slice of the generated noise. The ISC loss is:
\begin{equation}
  \mathcal{L}_{\text{ISC}}
  =-\frac{1}{D{-}1}\sum_{z=1}^{D-1}
    \bigl\|\bm{S}_{z+1}-\bm{S}_z\bigr\|_2.
  \label{eq:isc}
\end{equation}
Minimizing $\mathcal{L}_{\text{ISC}}$ (i.e., maximizing inter-slice spectral
divergence) forces drastic spectral shifts between adjacent slices,
injecting high-frequency ``spectral jitter'' along $z$ that
fundamentally contradicts the smooth anatomical continuity 3D
convolutions rely on.
We measure divergence in the \textit{frequency domain} rather than
the spatial domain because magnitude spectra are
translation-invariant~\cite{yin2019fourier,wang2020high}, providing a more stable
and texture-focused measure of inter-slice dissimilarity.
%------- L_SP -------
\subsubsection{Semantic Prediction Disruption.}
Effective data protection requires that perturbations corrupt not only
low-level pixel statistics but also the high-level semantic mapping
learned by the surrogate. If the surrogate can still produce similar
predictions on clean and perturbed inputs, the learned representations
remain intact and protection fails. We therefore explicitly maximize
the discrepancy between the surrogate's outputs on clean and perturbed
volumes, forcing perturbations to propagate through the entire network
and disrupt its semantic understanding.

Let $\bm{L}_{z}\!=\!\mathcal{F}_s(\bm{x};\theta)$ and
$\bm{L}_{z'}\!=\!\mathcal{F}_s(\bm{x}_p;\theta)$ denote the
clean and perturbed logit volumes
($\in\mathbb{R}^{D\times H\times W\times K}$), respectively.
We maximize the $\ell_1$ distance between them:
\begin{equation}
  \mathcal{L}_{\text{SPD}}
  =-\bigl\|\,\bm{L}_{z'} - \bm{L}_{z}\bigr\|_1.
  \label{eq:sp}
\end{equation}
Compared with the $\ell_2$ norm, the $\ell_1$ norm is less dominated by a few large-deviation voxels, thereby encouraging the generator to degrade predictions broadly across the volume rather than at isolated locations.

%------- Optimization -------
\subsubsection{Optimization.}
The total generator objective combines segmentation, logits
divergence, and inter-slice consistency disruption:
\begin{equation}
  \mathcal{L}_{\mathrm{total}}
  =\mathcal{L}_{\mathrm{seg}}^{\mathcal{'}}\!\big(\mathcal{F}_s(\bm{x}_p;\theta),\,\bm{y}\big)
  +\lambda_{\text{SPD}}\,\mathcal{L}_{\text{SPD}}
  +\lambda_{\text{ISC}}\,\mathcal{L}_{\text{ISC}},
  \label{eq:total}
\end{equation}
where $\mathcal{L}_{\text{seg}}^{'}$ (DiceCE) anchors training stability
by ensuring the surrogate continues to learn on perturbed data (thus
providing meaningful gradients), while $\mathcal{L}_{\text{SPD}}$
and $\mathcal{L}_{\text{ISC}}$ inject semantic and structural corruption,
respectively. $\lambda_{\text{SPD}}$ and $\lambda_{\text{ISC}}$ balance the three terms.

Following the bi-level framework of Eq.~\eqref{eq:ue}, the surrogate is frozen and the
generator $\mathcal{G}$ is updated over the full training set by
minimizing $\mathcal{L}_{\mathrm{total}}$:
\begin{equation}
  \min_{\theta_{\mathcal{F}_s}}\;
  \mathbb{E}_{(\bm{x},\bm{y})\sim\mathcal{D}_{\mathrm{clean}}}
  \!\Bigl[\min_{\theta_{\mathcal{G}}}
    \mathcal{L}_{\mathrm{total}}\Bigr].
  \label{eq:bilevel}
\end{equation}
As detailed in Algorithm~\ref{alg:VoxShield}, training alternates
between an S-step that updates the surrogate $\mathcal{F}_s$ on
perturbed data and an N-step that updates the generator $\mathcal{G}$
via $\mathcal{L}_{\mathrm{total}}$. After convergence, the trained
generator is applied to all clean volumes to produce
$\mathcal{D}_{\mathrm{protect}}$.

\begin{table}[t]
\centering
\small
\setlength{\tabcolsep}{3pt}
\begin{tabular}{p{0.95\linewidth}}
\toprule
\textbf{Algorithm 1:} VoxShield Training Procedure \\
\midrule
\textbf{Input:} Clean dataset $\mathcal{D}_{\mathrm{clean}}$,
  surrogate $\mathcal{F}_s$, generator $\mathcal{G}$, budget
  $\epsilon$, weights $\lambda_{\text{SPD}},\lambda_{\text{ISC}}$, learning rates
  $\alpha_s,\alpha_g$, S-step batches $T_s$ \\
\textbf{Output:} Protected dataset $\mathcal{D}_{\mathrm{protect}}$ \\[2pt]
1: \textbf{for} epoch $= 1$ \textbf{to} $n$ \textbf{do} \\
2: \quad \textit{// N-step: update generator} \\
3: \quad \textbf{for} each $(\bm{x},\bm{y})$ in
  $\mathcal{D}_{\mathrm{clean}}$ \textbf{do} \\
4: \quad\quad $\bm{\delta}\leftarrow\mathcal{G}(\bm{x})$;~~
  $\bm{x}_p\leftarrow\mathrm{Clamp}_{[0,1]}[\bm{x}+\bm{\delta}]$ \\
5: \quad\quad Compute $\mathcal{L}_{\mathrm{total}}$
  \hfill \textit{// Eq.~\eqref{eq:total}} \\
6: \quad\quad $\theta_{\mathcal{G}}\leftarrow
  \theta_{\mathcal{G}}-\alpha_g\nabla_{\theta_{\mathcal{G}}}
  \mathcal{L}_{\mathrm{total}}$ \\
7: \quad \textbf{end for} \\
8: \textbf{end for} \\
9: \textbf{for} each $(\bm{x},\bm{y})$ in
  $\mathcal{D}_{\mathrm{clean}}$ \textbf{do} \\
10: \quad $\mathcal{D}_{\mathrm{protect}}\leftarrow
  \mathcal{D}_{\mathrm{protect}}\cup
  \{(\mathrm{Clamp}_{[0,1]}[\bm{x}+\mathcal{G}(\bm{x})],
  \;\bm{y})\}$ \\
11: \textbf{end for} \\
12: \textbf{return} $\mathcal{D}_{\mathrm{protect}}$ \\
\bottomrule
\end{tabular}
\label{alg:VoxShield}
\end{table}

% ═══════════════════════════════════════════════════════════════════
%                  3  EXPERIMENTS AND RESULTS
% ═══════════════════════════════════════════════════════════════════
\section{Experiments and Results}

\subsection{Experimental Setup}

\textbf{Datasets.}
We evaluated the protection effectiveness of VoxShield on two 3D MIS datasets:
\textbf{BraTS19}~\cite{menze2014multimodal} (brain tumor MRI) and
\textbf{FLARE21}~\cite{ma2021abdomenct} (abdominal organ CT). These two datasets
cover different anatomical regions and imaging techniques, providing comprehensive
validation for the generalizability of our approach.


\textbf{Models.}
We used a 3D UNet~\cite{cciccek20163d} implemented in MONAI as the surrogate model
for UE generation. To evaluate cross-architecture transferability, we trained four
victim models: \textbf{3D-UNet}~\cite{cciccek20163d}, \textbf{Attention UNet}~\cite{oktay2018attention},
\textbf{UNet++}~\cite{zhou2018unetpp}, and \textbf{TransUNet}~\cite{chen2021transunet}.


\textbf{Baselines.}
We compared against: (1)~\textit{Clean}; (2)~\textit{Gaussian Noise};
(3)~\textit{EM}~\cite{huang2021unlearnable}; 
(4)~\textit{LSP}~\cite{yu2022availability}; (5)~\textit{SEP}~\cite{liu2023stable};
(6)~\textit{TAP}~\cite{fowl2021adversarial};
(7)~\textit{PUE}~\cite{wang2025provably}; 
(8)~\textit{UMed}~\cite{lin2024safeguarding}. All baselines were
adapted to 3D volumetric inputs.

\textbf{Metrics and Implementation.}
We reported the Dice Similarity Coefficient (DSC, \%) and 95th-percentile Hausdorff Distance (HD95, mm). Lower DSC and higher HD95 on victim models reflect more effective data protection. The perturbation budget was
$\epsilon=4/255$. We trained VoxShield for 100 epochs with Adam optimizer (lr$=$1e-4) on a single NVIDIA
RTX 5090 GPU. The loss weights were set to $\lambda_{\text{ISC}}=0.2$ and $\lambda_{\text{SPD}}=0.05$.

\subsection{Main Results}

% ── Table 1: Main comparison on both datasets ──
% ── Table 1: Main comparison on both datasets ──
\begin{table}[t]
\centering
\caption{Protection effectiveness on victim 3D\,UNet.
  DSC\,(\%) $\downarrow$ and HD95\,(mm) $\uparrow$ indicate stronger
  protection; PSNR\,(dB) $\uparrow$ and SSIM $\uparrow$ measure
  perturbation imperceptibility.
  Best UE results in \textbf{bold}.}\label{tab:main}
\setlength{\tabcolsep}{3pt}
\footnotesize
\begin{tabular}{l cccc cccc}
\toprule
\multirow{2}{*}{Method}
  & \multicolumn{4}{c}{BraTS19}
  & \multicolumn{4}{c}{FLARE21} \\
\cmidrule(lr){2-5}\cmidrule(lr){6-9}
  & DSC$\downarrow$ & HD95$\uparrow$
  & PSNR$\uparrow$ & SSIM$\uparrow$
  & DSC$\downarrow$ & HD95$\uparrow$
  & PSNR$\uparrow$ & SSIM$\uparrow$ \\
\midrule
Clean
  & 80.0 &  16.0 & --   & --
  & 88.6 &  4.6 & --   & -- \\
Gaussian Noise
  & 81.4 &  14.2 & 44.4   & 0.9688
  & 89.2 &  5.2 & \textbf{50.5}   & 0.9943 \\
EM~\cite{huang2021unlearnable}
  & 69.3 & 30.2 & 39.7   & 0.9219
  & 57.1 & 4.14 & 30.0   & 0.9312 \\
LSP~\cite{yu2022availability}
  & 80.5 & 15.0 & 39.6   & 0.9916
  & 88.5 & 7.6 & 43.0   & 0.9970 \\
SEP~\cite{liu2023stable}
  & 72.3 &  16.7 & 33.9   & 0.7544
  & 88.4 &  7.0 & 40.8   & 0.9511 \\
TAP~\cite{liu2023stable}
  & 68.4 &  22.7 & 38.0   & 0.8792
  & 45.7 &  40.9 & 37.7   & 0.9084 \\
PUE~\cite{wang2025provably}
  & 60.0 &  26.5 & 40.2   & 0.9279
  & 36.1 &  45.4 & 37.9   & 0.9155 \\
UMed~\cite{lin2024safeguarding}
  & 40.2 &  34.7 & 54.9   & 0.9990
  & 26.3 &  49.7 & 48.0   & 0.9973 \\
\midrule
\textbf{Our VoxShield}
  & \textbf{$\approx$\,0.0} & \textbf{217.4} & \textbf{56.3} & \textbf{0.9994}
  & \textbf{6.8} & \textbf{76.0}    & 49.0 & \textbf{0.9980} \\
\bottomrule
\end{tabular}
\end{table}

Table~\ref{tab:main} reported the protection performance on both datasets. VoxShield achieved a dramatic
reduction in victim model DSC (from 89\% to 7\% on FLARE21), substantially outperforming all baselines
on both datasets. 2D-originated methods (EM, SEP, LSP, PUE, UMed) showed limited effectiveness on 3D
data---their slice-independent noise patterns were partially smoothed by 3D convolution kernels,
confirming our hypothesis that disrupting inter-slice consistency is essential for effective 3D protection.

\textbf{Cross-Architecture Transferability}

% ── Table 3: Transfer ──
\begin{table}[t]
\centering
\caption{Cross-architecture transferability on FLARE21 (Mean DSC\,\%). Noise generated with UNet surrogate, evaluated on different victim architectures.}\label{tab:transfer}
\setlength{\tabcolsep}{6pt}
\begin{tabular}{lcccc}
\toprule
Method & 3D-UNet & Att-UNet & UNet++ & TransUNet \\
\midrule
Clean     & 87.6 & 83.2   & 83.4   & 80.7 \\
EM       & 69.3   & 48.6   & 68.6   & 74.2 \\
\textbf{VoxShield}  & \textbf{6.8} & \textbf{9.0} & \textbf{10.8} & \textbf{3.1} \\
\bottomrule
\end{tabular}
\end{table}

Table~\ref{tab:transfer} demonstrated that VoxShield perturbations transferred effectively across
architectures. Since the perturbations targeted the fundamental volumetric continuity prior shared by
\textit{all} 3D architectures (not architecture-specific features), the protection generalized from the
UNet surrogate to Attention UNet, UNet++, and even Transformer-based TransUNet.

\textbf{Ablation Study}

\begin{table}[t]
\centering
\caption{Ablation study on FLARE21. ``Per-slice Gen.'' generates
  perturbations independently for each slice via the noise generator,
  without explicitly disrupting inter-slice consistency.
  Best results in \textbf{bold}.}\label{tab:ablation}
\setlength{\tabcolsep}{6pt}
\begin{tabular}{lcc}
\toprule
Configuration & DSC\,(\%)\,$\downarrow$ & HD95\,(mm)\,$\uparrow$ \\
\midrule
\textbf{VoxShield (full)}
  & \textbf{6.8}  & \textbf{76.0} \\
\midrule
w/o $\mathcal{L}_{\text{SPD}}$ (no logits divergence)
  & 24.8    & 53.4 \\
w/o $\mathcal{L}_{\text{ISC}}$ (no ISC disruption)
  & 32.9  & 41.2 \\
Per-slice Gen.\ (no cross-slice disruption)
  & 74.9    & 31.4 \\
\bottomrule
\end{tabular}
\end{table}


Table~\ref{tab:ablation} validated the contribution of each component on FLARE21. Removing
$\mathcal{L}_{\text{ISC}}$ caused the largest performance drop, confirming that inter-slice spectral
diversity is the most critical attack vector against 3D segmentation. Removing $\mathcal{L}_{\text{SPD}}$
also degraded protection, as the generator lost its ability to corrupt deep semantic representations.

\textbf{Inter-Slice Frequency Consistency Disruption.}
We visualized the noise patterns across consecutive $z$-slices for different methods.
VoxShield produced spectrally diverse perturbations that shifted drastically between
adjacent slices, while 2D methods (EM, SEP, LSP, PUE, UMed) generated temporally coherent noise that
3D kernels could readily filter. Quantitatively, the mean inter-slice spectral
correlation of VoxShield noise was significantly lower than all baselines, directly
evidencing the effectiveness of $\mathcal{L}_{\text{ISC}}$.
Fig.~\ref{fig:isc} illustrates these differences.

\begin{figure}[t]
\centering
\includegraphics[width=\textwidth]{fig/result.pdf}
\caption{Inter-slice frequency consistency disruption. Each column shows two consecutive $z$-slices iamge. 2D methods (EM, LSP)
produce spectrally consistent noise across slices (high inter-slice correlation), which 3D convolutional kernels can easily filter out. In contrast, VoxShield generates spectrally diverse perturbations with significantly lower inter-slice spectral correlation, effectively disrupting the frequency consistency that 3D models rely on.}
\label{fig:isc}
\end{figure}


\textbf{Segmentation Degradation Patterns.}
Victim model predictions on VoxShield-protected data exhibited a characteristic 3D-specific degradation:
segmentation masks became \textit{inter-slice discontinuous}---fragmented, with boundaries jumping between
slices. In contrast, 2D UE methods produced only mild global accuracy reduction while retaining smooth
inter-slice predictions. This degradation pattern directly demonstrated that VoxShield successfully
disrupted the volumetric continuity that 3D models relied upon.


% ═══════════════════════════════════════════════════════════════════
%                       4  CONCLUSION
% ═══════════════════════════════════════════════════════════════════
\section{Conclusion}

We present VoxShield, the first unlearnable example framework specifically designed for 3D medical image segmentation. By identifying inter-slice anatomical consistency as the critical---and exploitable---prior of 3D segmentation models, we design frequency-domain objectives that disrupt this prior at both the noise-structure level ($\mathcal{L}_{\text{ISC}}$) and the semantic-prediction level ($\mathcal{L}_{\text{SPD}}$). VoxShield achieves state-of-the-art protection across two benchmarks and four victim architectures, demonstrating that effective 3D data protection requires attacking the volumetric priors that 3D models fundamentally depend upon.

% ═══════════════════════════════════════════════════════════════════
%                       BIBLIOGRAPHY
% ═══════════════════════════════════════════════════════════════════
%\bibliographystyle{splncs04}
%\bibliography{references}

\begin{thebibliography}{20}

\bibitem{huang2021unlearnable}
Huang, H., Ma, X., Erfani, S.M., Bailey, J., Wang, Y.:
Unlearnable examples: Making personal data unexploitable.
In: ICLR (2021)

\bibitem{tcia}
The Cancer Imaging Archive (TCIA),
\url{https://www.cancerimagingarchive.net/}, last accessed 2025/01/15

\bibitem{pubmed}
PubMed, \url{https://pubmed.ncbi.nlm.nih.gov/}, last accessed 2025/01/15


\bibitem{gdpr2016}
{European Parliament and Council of the European Union}:
Regulation ({EU}) 2016/679 (General Data Protection Regulation).
Official Journal of the European Union, L~119, 1--88 (2016)

\bibitem{fowl2021adversarial}
Fowl, L., Goldblum, M., Chiang, P.Y., Geiping, J., Czaja, W., Goldstein, T.:
Adversarial examples make strong poisons.
In: NeurIPS (2021)

\bibitem{liu2023stable}
Liu, Y., Xu, K., Chen, X., Sun, L.:
Stable unlearnable example: Enhancing the robustness of unlearnable examples via stable error-minimizing noise.
In: AAAI, pp.~3783--3791 (2024)

\bibitem{yu2022availability}
Yu, D., Zhang, H., Chen, W., Yin, J., Liu, T.Y.:
Availability attacks create shortcuts.
In: KDD, pp.~2367--2376 (2022)


\bibitem{cciccek20163d}
\c{C}i\c{c}ek, \"O., Abdulkadir, A., Lienkamp, S.S., Brox, T., Ronneberger, O.:
3D U-Net: Learning dense volumetric segmentation from sparse annotation.
In: MICCAI, pp.~424--432 (2016)

\bibitem{milletari2016v}
Milletari, F., Navab, N., Ahmadi, S.A.:
V-Net: Fully convolutional neural networks for volumetric medical image segmentation.
In: 3DV, pp.~565--571 (2016)

\bibitem{isensee2021nnu}
Isensee, F., Jaeger, P.F., Kohl, S.A., Petersen, J., Maier-Hein, K.H.:
nnU-Net: A self-configuring method for deep learning-based biomedical image segmentation.
Nature Methods \textbf{18}(2), 203--211 (2021)

\bibitem{menze2014multimodal}
Menze, B.H., Jakab, A., Bauer, S., et~al.:
The multimodal brain tumor image segmentation benchmark (BRATS).
IEEE TMI \textbf{34}(10), 1993--2024 (2015)

\bibitem{ma2021abdomenct}
Ma, J., Zhang, Y., Gu, S., et~al.:
AbdomenCT-1K: Is abdominal organ segmentation a solved problem?
IEEE TPAMI \textbf{44}(10), 6695--6714 (2022)

\bibitem{oktay2018attention}
Oktay, O., Schlemper, J., Folgoc, L.L., et~al.:
Attention U-Net: Learning where to look for the pancreas.
In: MIDL (2018)

\bibitem{zhou2018unetpp}
Zhou, Z., Siddiquee, M.M.R., Tajbakhsh, N., Liang, J.:
UNet++: A nested U-Net architecture for medical image segmentation.
In: DLMIA/ML-CDS@MICCAI, pp.~3--11 (2018)

\bibitem{chen2021transunet}
Chen, J., Lu, Y., Yu, Q., et~al.:
TransUNet: Transformers make strong encoders for medical image segmentation.
arXiv preprint arXiv:2102.04306 (2021)

\bibitem{hatamizadeh2022unetr}
Hatamizadeh, A., Tang, Y., Nath, V., Yang, D., Myronenko, A., Landman, B., Roth, H.R., Xu, D.:
UNETR: Transformers for 3D medical image segmentation.
In: WACV, pp.~574--584 (2022)

\bibitem{tang2022swinunetr}
Tang, Y., Yang, D., Li, W., Roth, H.R., Landman, B., Xu, D., Nath, V., Hatamizadeh, A.:
Self-supervised pre-training of swin transformers for 3D medical image analysis.
In: CVPR, pp.~20730--20740 (2022)

\bibitem{kaissis2021end}
Kaissis, G.A., Makowski, M.R., R\"{u}ckert, D., Braren, R.F.:
Secure, privacy-preserving and federated machine learning in medical imaging.
Nature Machine Intelligence \textbf{2}(6), 305--311 (2020)

\bibitem{fu2022robust}
Fu, S., He, F., Liu, Y., Shen, L., Tao, D.:
Robust unlearnable examples: Protecting data against adversarial learning.
In: ICLR (2022)

\bibitem{sadasivan2023cuda}
Sadasivan, V.S., Soltanolkotabi, M., Feizi, S.:
CUDA: Convolution-based unlearnable datasets.
In: CVPR, pp.~3862--3871 (2023)

\bibitem{sun2024unseg}
Sun, J., Chen, Q., Luo, J., et al.:
UnSeg: One universal unlearnable example generator is enough against all image segmentation.
In: NeurIPS (2024)

\bibitem{wang2024unlearnable3d}
Wang, X., Zeng, Z., Liu, Y., et al.:
Unlearnable 3D point clouds: Class-wise transformation is all you need.
In: NeurIPS (2024)

\bibitem{dong2022mnet}
Dong, Z., He, Y., Qi, X., Chen, Y., Shu, H., Coatrieux, J.L., Yang, G., Li, S.:
MNet: Rethinking 2D/3D networks for anisotropic medical image segmentation.
In: IJCAI, pp.~870--876 (2022)


\bibitem{yin2019fourier}
Yin, D., Lopes, R.G., Shlens, J., Cubuk, E.D., Gilmer, J.:
A Fourier perspective on model robustness in computer vision.
In: NeurIPS, pp.~13255--13265 (2019)

\bibitem{wang2020high}
Wang, H., Wu, X., Huang, Z., Xing, E.P.:
High-frequency component helps explain the generalization of convolutional neural networks.
In: CVPR, pp.~8684--8694 (2020)

\bibitem{lin2024safeguarding}
Lin, X., Yu, Y., Xia, S., Jiang, J., Wang, H., Yu, Z., Liu, Y., Fu, Y., Wang, S., Tang, W., Kot, A.:
Safeguarding medical image segmentation datasets against unauthorized training via contour- and texture-aware perturbations.
arXiv preprint arXiv:2403.14250 (2024)

\bibitem{wang2025provably}
Wang, D., Xue, M., Li, B., Camtepe, S., Zhu, L.:
Provably unlearnable data examples.
In: Network and Distributed System Security Symposium (NDSS) (2025)

\end{thebibliography}

\end{document}